\newpage
\begin{center}
  \parbox{0.95\textwidth}{\centering\textbf{\large 3. ВНЕДРЕНИЕ СИСТЕМЫ И ЭКСПЕРИМЕНТАЛЬНОЕ ИССЛЕДОВАНИЕ}}
\end{center}
\refstepcounter{chapter}
\addcontentsline{toc}{chapter}{3. ВНЕДРЕНИЕ СИСТЕМЫ И ЭКСПЕРИМЕНТАЛЬНОЕ ИССЛЕДОВАНИЕ}

\section{Цель экспериментальной части}

Экспериментальная часть проверяет, насколько устойчиво можно восстановить
структуру случайного бинарного изображения $6 \times 6$ по EEG-сигналу,
зарегистрированному во время визуального воображения. Один экспериментальный
пример соответствует одной эпохе воображения и одному целевому изображению.
Цель представляется как многометочная задача с 36 бинарными выходами:

\begin{equation}
\mathbf{y}, \hat{\mathbf{y}} \in \{0,1\}^{36}.
\end{equation}

Основной корпус экспериментальных строк содержит 180 эпох случайного
визуального воображения. Каждая строка сохраняет идентификатор испытуемого,
номер экспериментальной пробы и ключ целевого блока, что позволяет явно
контролировать утечки между обучением и оценкой. Входом моделей служат только
EEG-каналы; EOG-записи используются как часть исходного экспериментального
контекста, но в признаки и нейросетевые входы данной главы не включаются.

\section{Программная реализация}

Пайплайн реализован на языке Python и разделён на этапы загрузки данных,
построения признаков, спектральных преобразований, обучения классических и
нейросетевых моделей, а также оценки сохранённых артефактов. Для численных
операций используются NumPy и SciPy, для классических моделей -- scikit-learn,
для нейросетевых моделей -- PyTorch.

Экспериментальные запуски сохраняются как неизменяемые директории артефактов:
в них входят конфигурация, описание разбиения, аудит утечек, предсказания,
оценочные метрики, bootstrap-интервалы, сведения о среде исполнения и
контрольная инвентаризация файлов. Итоговое сравнение загружает только
метаданные и массивы предсказаний; сериализованные scikit-learn pipeline и
PyTorch checkpoint-файлы не десериализуются при безопасной проверке результатов.

\begin{figure}[!t]
  \centering
  \includegraphics[width=150mm]{experiment_pipeline.jpg}
  \caption{Общая схема экспериментального пайплайна реконструкции случайного изображения.}
  \label{fig:experiment-pipeline}
\end{figure}

\section{Подготовка входных представлений}

Используются уже предобработанные EEG-записи, процедура исходной
предобработки описана в главе 2. Для экспериментов из каждой эпохи выбирается
полуоткрытый интервал $[0.5, 15.5)$ с, после чего сигнал приводится к
частоте анализа 125 Гц. Целевая матрица изображения векторизуется построчно
в вектор длины 36.

Для классических моделей строятся признаки полного 15-секундного окна:
временные статистики, спектральные мощности, пространственные ковариационные
представления и локальные паттерны. При обучении все операции выбора группы
признаков, отбора признаков, масштабирования и настройки гиперпараметров
выполняются только внутри обучающих частей соответствующего протокола.
Полное окно рассматривается как одна строка модели, поскольку оно соответствует
одной размеченной эпохе воспоминания и одному целевому изображению; зависимое
оконное дробление без отдельной схемы группировки не используется в итоговом
сравнении.

Для нейросетевых моделей строятся спектральные входы четырёх типов:
FFT/Welch, Morlet, Superlet и STFT. FFT/Welch даёт тензор вида
канал--частота, а Morlet, Superlet и STFT -- тензор канал--частота--время.
Перед обучением выполняется логарифмирование мощности и нормализация по
частотным bin-ам, параметры которой оцениваются только по обучающей части.
Тестовые спектральные массивы не используются при подборе нормировки,
длительности обучения или весов классов.

Сводка входных представлений приведена в табл.
\ref{tab:experiment-input-representations}. Все представления строятся из
одного и того же 15-секундного EEG-фрагмента, поэтому различия между
моделями не связаны с использованием разных временных интервалов.

\begin{table}[!t]
  \centering
  \small
  \setlength{\tabcolsep}{3pt}
  \renewcommand{\arraystretch}{1.15}
  \caption{Входные представления для экспериментальных моделей.}
  \label{tab:experiment-input-representations}
  \begin{tabularx}{\textwidth}{|L{36mm}|L{39mm}|Y|}
    \hline
    Представление & Используемые модели & Назначение \\
    \hline
    Полное временное окно &
    Logistic Regression и дополнительные классические модели &
    Признаки времени, спектральных мощностей, пространственных ковариаций и
    локальных паттернов; используются для табличных scikit-learn моделей. \\
    \hline
    FFT/Welch &
    EEGNet, DeepConvNet, ShallowConvNet &
    Тензор канал--частота; используется как компактный спектральный вход без
    временной оси. \\
    \hline
    Morlet, Superlet, STFT &
    EEGNet, DeepConvNet, ShallowConvNet &
    Тензоры канал--частота--время; используются для проверки CNN-моделей на
    time-frequency представлениях. \\
    \hline
  \end{tabularx}
\end{table}

\section{Модели}

В сравнении участвуют три группы методов. Первая группа -- опорная
логистическая регрессия: для каждого пикселя обучается отдельная L1 / L2
Logistic Regression модель, а группа признаков и гиперпараметры выбираются
по внутренней кросс-валидации с группировкой по испытуемым. В опубликованном
межсубъектном запуске отбор признаковой группы только по обучающей части
выбрал локально-паттерновую группу 1D-LBP.

Вторая группа -- дополнительные классические модели: Linear SVM, Ridge Classifier,
Ridge Regression, ElasticNet/Lasso, Random Forest Regressor и PLS Regression
в независимых или многовыходных вариантах. Для регрессионных моделей выходные
значения интерпретируются как непрерывные оценки, обрезанные к диапазону
$[0,1]$ перед пороговым решением; они не рассматриваются как
калиброванные вероятности.

Третья группа -- 12 нейросетевых вариантов, полученных декартовым
произведением трёх архитектур EEGNet, DeepConvNet и ShallowConvNet и четырёх
спектральных представлений FFT/Welch, Morlet, Superlet и STFT. Каждая модель
имеет 36 логитов на выходе. Финальные нейросетевые предсказания формируются
ансамблем трёх запусков с начальными зёрнами 42, 43 и 44, после чего
применяется фиксированный порог 0.5.

Перечень семейств моделей и их роль в сравнении приведены в табл.
\ref{tab:experiment-model-families}. Таблица фиксирует не только названия
моделей, но и то, какой вопрос закрывает каждая группа.

\begin{table}[!t]
  \centering
  \small
  \setlength{\tabcolsep}{3pt}
  \renewcommand{\arraystretch}{1.15}
  \caption{Нейросетевые и классические модели в экспериментальном пайплайне.}
  \label{tab:experiment-model-families}
  \begin{tabularx}{\textwidth}{|L{34mm}|L{30mm}|Y|}
    \hline
    Семейство & Количество вариантов & Роль в эксперименте \\
    \hline
    Logistic Regression &
    1 референсный запуск &
    Основная точка сравнения: независимая бинарная модель для каждого
    пикселя, подбор признаков и гиперпараметров с группировкой по испытуемым внутри
    обучающей части. \\
    \hline
    Дополнительные классические модели &
    9 обучаемых не-референсных вариантов &
    Проверяют, дают ли линейные, регуляризованные, лесные и PLS-модели
    устойчивое улучшение над референсом на тех же табличных признаках. \\
    \hline
    Torch CNN-модели &
    12 вариантов &
    Проверяют спектральные адаптации EEGNet, DeepConvNet и ShallowConvNet на
    FFT/Welch, Morlet, Superlet и STFT входах; каждый вариант оценивается как
    ансамбль трёх seed-запусков. \\
    \hline
    Базовые модели без EEG-сигнала &
    3 стратегии &
    Задают нижний контекст для интерпретации метрик: majority, попиксельная
    частота и бернуллиевская стратегия с фиксированным seed без использования
    EEG-сигнала. \\
    \hline
  \end{tabularx}
\end{table}

\section{Протоколы оценки}

Используются два протокола, которые отвечают на разные вопросы и поэтому не
объединяются в один общий результат. Их параметры приведены в табл.
\ref{tab:experiment-protocols}.

\begin{table}[!t]
  \centering
  \small
  \setlength{\tabcolsep}{3pt}
  \renewcommand{\arraystretch}{1.15}
  \caption{Протоколы экспериментальной оценки.}
  \label{tab:experiment-protocols}
  \begin{tabularx}{\textwidth}{|L{33mm}|L{31mm}|L{31mm}|Y|}
    \hline
    Протокол & Обучение & Тест & Контролируемая граница \\
    \hline
    Межсубъектный &
    141 строка, 26 испытуемых &
    39 строк, 7 испытуемых &
    Испытуемые, ключи наблюдений, случайные зёрна и целевые изображения не
    пересекаются между обучением и тестом. \\
    \hline
    Двунаправленный межпробный &
    81 строка в каждом направлении &
    81 строка в каждом направлении &
    Испытуемые совпадают между обучением и тестом, но номера проб, ключи
    наблюдений, случайные зёрна и целевые изображения не пересекаются. В итоговой
    сводке объединяются два направления, всего 162 тестовые строки от 27
    испытуемых. \\
    \hline
  \end{tabularx}
\end{table}

В межсубъектном протоколе оценивается перенос на новых испытуемых. В
двунаправленном межпробном протоколе оценивается перенос между двумя
экспериментальными пробами внутри тех испытуемых, для которых доступны обе
пробы; испытуемые 14, 24, 27, 28, 29 и 32 исключаются из этого протокола
из-за отсутствия полной пары проб. Для обеих схем аудит артефактов подтвердил
отсутствие запрещённых пересечений по ключам наблюдений, случайным зёрнам и полным
целевым изображениям.

\section{Метрики и статистическая проверка}

Основная метрика -- средняя попиксельная сбалансированная точность
(per-pixel balanced accuracy) по 36 бинарным задачам. Она выбрана вместо
простой попиксельной точности (pixel accuracy), поскольку число активных
и неактивных пикселей в конкретных тестовых подмножествах может отличаться.
Дополнительно вычисляются битовая точность, macro F1, micro IoU, Hamming loss,
Brier score или score-MSE, а также точность полного совпадения для полной
матрицы $6 \times 6$.

Базовые модели без EEG-сигнала включают глобальное majority-предсказание,
попиксельную частотную стратегию и бернуллиевскую стратегию с фиксированным
seed. Для неопределённости используется субъектно-кластерный bootstrap: в
межсубъектном протоколе ресэмплируются тестовые испытуемые, а в объединённой
межпробной оценке -- целые испытуемые с их удержанными строками. Для итогового сравнения
используются 2,000 принятых bootstrap-выборок на протокол.

Сравнение моделей с Logistic Regression выполняется как парное
bootstrap-сравнение по сбалансированной точности. Поправка Holm применяется отдельно внутри каждого
протокола по 21 не-референсной обучаемой модели. Поэтому описательные ранги
моделей в этой главе не интерпретируются как доказательство статистически
надёжного превосходства.

\section{Результаты межсубъектного протокола}

Основные результаты межсубъектного протокола приведены в табл.
\ref{tab:cross-subject-results}. Все значения получены из сохранённых
артефактов оценки.

\begin{table}[!t]
  \centering
  \small
  \setlength{\tabcolsep}{3pt}
  \renewcommand{\arraystretch}{1.15}
  \caption{Выбранные результаты межсубъектного протокола.}
  \label{tab:cross-subject-results}
  \begin{tabularx}{\textwidth}{|L{43mm}|L{24mm}|L{22mm}|L{18mm}|L{18mm}|Y|}
    \hline
    Модель & Balanced accuracy & 95\% bootstrap CI & Битовая точность & Micro IoU & Комментарий \\
    \hline
    Logistic Regression &
    0.509991 & [0.496384; 0.521077] & 0.514245 & 0.334634 &
    Референсная модель сравнения; точность полного совпадения равна 0. \\
    \hline
    Ridge Regression independent &
    0.518382 & [0.486733; 0.561786] & 0.518519 & 0.326693 &
    Описательный лидер межсубъектного сравнения среди обучаемых моделей. \\
    \hline
    ShallowConvNet Morlet &
    0.513443 & [0.488845; 0.536044] & 0.524929 & 0.330321 &
    Лучший межсубъектный результат среди Torch-вариантов. \\
    \hline
  \end{tabularx}
\end{table}

\begin{figure}[!t]
  \centering
  \includegraphics[width=150mm]{experiment_cross_subject_ranking.pdf}
  \caption{Ранжирование моделей по balanced accuracy в межсубъектном протоколе. Пунктирная линия
  соответствует уровню 0.5.}
  \label{fig:cross-subject-ranking}
\end{figure}

\begin{figure}[!t]
  \centering
  \includegraphics[width=150mm]{experiment_cross_subject_delta_logistic.pdf}
  \caption{Парные bootstrap-контрасты balanced accuracy относительно Logistic Regression в
  межсубъектном протоколе. Интервалы являются поточечными bootstrap-интервалами.}
  \label{fig:cross-subject-delta-logistic}
\end{figure}

Результаты находятся близко к случайному уровню balanced accuracy 0.5.
Описательный лидер Ridge Regression independent превосходит Logistic
Regression по точечной оценке примерно на 0.0084, однако итоговая
multiplicity-aware проверка не позволяет объявлять это надёжным
превосходством. Минимальное bootstrap p-value после поправки Holm по финальному
сравнению равно 0.273000.

\section{Результаты двунаправленного межпробного протокола}

Внутри-субъектный протокол имеет другую интерпретацию: он допускает совпадение
испытуемых между обучением и тестом, но проверяет перенос между различными
пробами. Два направления обучаются и оцениваются независимо, после чего
объединяются только на уровне итогового сравнения.

Описательным лидером объединённого двунаправленного межпробного сравнения является
DeepConvNet STFT: средняя balanced accuracy равна 0.512011. Для направления
проба 2 $\rightarrow$ проба 1 сохранённый артефакт даёт balanced accuracy
0.513886. 95\% интервал: [0.499969; 0.526317], битовая точность --
0.517490, а micro IoU -- 0.364499.
В другом направлении, проба 1 $\rightarrow$ проба 2, balanced accuracy равна
0.509210. В обоих направлениях точность полного совпадения равна 0.

\begin{table}[!t]
  \centering
  \small
  \setlength{\tabcolsep}{3pt}
  \renewcommand{\arraystretch}{1.15}
  \caption{Выбранный результат двунаправленного межпробного протокола для DeepConvNet STFT.}
  \label{tab:within-subject-results}
  \begin{tabularx}{\textwidth}{|L{33mm}|L{24mm}|L{26mm}|L{18mm}|L{18mm}|Y|}
    \hline
    Оценка & Balanced accuracy & 95\% bootstrap CI & Битовая точность & Micro IoU & Комментарий \\
    \hline
    Проба 1 $\rightarrow$ проба 2 &
    0.509210 & [0.496211; 0.521737] & 0.512003 & 0.322381 &
    Одно из двух направлений внутри-субъектного переноса; точность полного совпадения равна 0. \\
    \hline
    Проба 2 $\rightarrow$ проба 1 &
    0.513886 & [0.499969; 0.526317] & 0.517490 & 0.364499 &
    Направление с большей точечной оценкой balanced accuracy; точность полного совпадения
    равен 0. \\
    \hline
    Объединённая двунаправленная оценка &
    0.512011 & [0.500668; 0.520872] & -- & -- &
    Объединённая сводка по 162 удержанным строкам от 27 испытуемых; используется
    только для итогового парного bootstrap-сравнения. \\
    \hline
  \end{tabularx}
\end{table}

\begin{figure}[!t]
  \centering
  \includegraphics[width=150mm]{experiment_within_subject_ranking.pdf}
  \caption{Ранжирование моделей по balanced accuracy в двунаправленном межпробном протоколе.}
  \label{fig:within-subject-ranking}
\end{figure}

\begin{figure}[!t]
  \centering
  \includegraphics[width=150mm]{experiment_within_subject_delta_logistic.pdf}
  \caption{Парные bootstrap-контрасты balanced accuracy относительно Logistic Regression в
  двунаправленном межпробном протоколе.}
  \label{fig:within-subject-delta-logistic}
\end{figure}

Таким образом, внутри-субъектный перенос между пробами также остаётся близким
к 0.5 по balanced accuracy. Лучшие Torch STFT варианты дают небольшие
положительные точечные отклонения, но эти отклонения рассматриваются как
описательные и протокол-специфичные.

\section{Примеры реконструкций}

Табл. \ref{tab:reconstruction-examples-real} показывает три реальные
тестовые реконструкции из межсубъектного артефакта ShallowConvNet Morlet
среди 39 тестовых примеров данного запуска. Строки выбраны по расстоянию
Хэмминга между истинными и предсказанными бинарными матрицами: минимальная,
медианная и максимальная ошибка.

\begin{table}[!t]
  \centering
  \scriptsize
  \setlength{\tabcolsep}{2pt}
  \renewcommand{\arraystretch}{1.15}
  \caption{Примеры реальных реконструкций на тестовых примерах.}
  \label{tab:reconstruction-examples-real}
  \begin{tabularx}{\textwidth}{|L{19mm}|L{19mm}|Y|Y|}
    \hline
    Пример & Ошибки & Истинная матрица $6 \times 6$ & Предсказанная матрица $6 \times 6$ \\
    \hline
    Минимальная ошибка &
    10 из 36 &
    \texttt{000000 / 001111 / 111110 / 000100 / 010011 / 011000} &
    \texttt{010000 / 001001 / 111100 / 100100 / 001110 / 011010} \\
    \hline
    Медианная ошибка &
    17 из 36 &
    \texttt{111110 / 101000 / 101000 / 100111 / 100101 / 011110} &
    \texttt{111000 / 101001 / 011000 / 101100 / 011010 / 011001} \\
    \hline
    Максимальная ошибка &
    24 из 36 &
    \texttt{010110 / 001111 / 001000 / 111101 / 000100 / 101011} &
    \texttt{111001 / 000001 / 110100 / 100100 / 110010 / 010000} \\
    \hline
  \end{tabularx}
\end{table}

\begin{figure}[!t]
  \centering
  \includegraphics[width=110mm]{experiment_reconstruction_examples_real.pdf}
  \caption{Графическое представление трёх реальных тестовых реконструкций: минимальная, медианная
  и максимальная ошибка Хэмминга среди выбранного межсубъектного запуска.}
  \label{fig:reconstruction-examples-real}
\end{figure}

\section{Интерпретация результатов}

Главный результат экспериментальной части состоит в строгой проверке задачи
реконструкции при контролируемых границах обучения и оценки. Завершённый
пайплайн воспроизводимо строит признаки и спектральные входы, фиксирует
разбиения, обучает классические и нейросетевые модели, сохраняет неизменяемые
артефакты, проводит bootstrap-оценку и выполняет безопасную проверку
сохранённых предсказаний. Внутри этой процедуры не получено убедительного
статистического свидетельства, что какая-либо обучаемая модель надёжно
превосходит Logistic Regression на текущей задаче реконструкции случайного
визуального образа.

Точность полного совпадения равна 0 для всех обучаемых моделей в финальном сравнении.
Это ожидаемо для задачи полной реконструкции 36 бинарных пикселей при малом
числе тестовых примеров и попиксельных метриках, близких к случайному уровню. Поэтому exact
метрика полного совпадения используется только как строгая вспомогательная метрика, а не как
основной критерий качества.

Отдельно важно, что битовая точность может быть немного выше 0.5 даже для
частотных или majority-стратегий из-за дисбаланса конкретных
пикселей. Поэтому выводы главы опираются прежде всего на balanced accuracy,
bootstrap-интервалы и парное сравнение с референсной Logistic Regression.

\section{Внешняя проверка пайплайна на открытых EEG-корпусах}

Для дополнительной проверки инженерной корректности пайплайна были выполнены
контрольные эксперименты на двух открытых корпусах из семейства BNCI. Эти
эксперименты не являются внешней валидацией самой задачи реконструкции
случайного изображения $6 \times 6$: целевые переменные, протоколы стимуляции
и длительности эпох отличаются от основного корпуса. Их роль другая -- проверить,
что реализованные компоненты загрузки, субъектно-раздельной оценки, построения
признаков, обучения моделей, сохранения артефактов и безопасной сводки
результатов работают на независимых EEG-задачах с существенно большим числом
эпох.

Первый внешний корпус, BNCI2014-001, относится к моторному воображению и
содержит 9 испытуемых, 22 EEG-канала и четыре класса: left hand, right hand,
feet и tongue. В итоговом leave-one-subject-out протоколе использовались
5184 эпохи. В сравнение вошли классический CSP+LDA baseline, Logistic
Regression на признаках текущего проекта и спектральные Torch-варианты
EEGNet, DeepConvNet и ShallowConvNet. Лучший результат показал CSP+LDA:
средняя balanced accuracy по 9 фолдам составила 0.385224 при стандартном
отклонении 0.108631. Это выше случайного уровня 0.25 для четырёх классов,
однако текущая адаптация нейросетевых моделей оказалась слабее: лучший Torch
вариант DeepConvNet STFT дал 0.308063.

Второй внешний корпус, BNCI2014-009, относится к P300 и содержит 10
испытуемых, 16 EEG-каналов и бинарную задачу Target/NonTarget. В
leave-one-subject-out протоколе использовались 17280 постстимульных эпох; в
каждом тестовом фолде было 1728 эпох, включая 288 Target и 1440 NonTarget
примеров. Лучший результат показала Logistic Regression на ERP/xDAWN-признаках:
средняя balanced accuracy составила 0.764062, macro F1 -- 0.711260, ROC AUC --
0.850690. Лучший raw Torch вариант, DeepConvNet на сырых ERP-эпохах, достиг
balanced accuracy 0.616146, а лучший спектральный Torch вариант, EEGNet FFT,
достиг 0.552674.

\begin{table}[!t]
  \centering
  \small
  \setlength{\tabcolsep}{3pt}
  \renewcommand{\arraystretch}{1.15}
  \caption{Контрольные внешние эксперименты на открытых EEG-корпусах.}
  \label{tab:external-benchmark-results}
  \begin{tabularx}{\textwidth}{|L{29mm}|L{34mm}|L{16mm}|L{32mm}|Y|}
    \hline
    Корпус & Задача и протокол & Число эпох & Лучший результат & Интерпретация \\
    \hline
    BNCI2014-001 &
    4 класса моторного воображения, leave-one-subject-out &
    5184 &
    CSP+LDA, balanced accuracy 0.385224 &
    Результат выше случайного уровня 0.25; проверены загрузка, LOSO-разбиение,
    признаки проекта и спектральные Torch-входы на независимом корпусе. \\
    \hline
    BNCI2014-009 &
    Target/NonTarget P300, leave-one-subject-out &
    17280 &
    ERP/xDAWN Logistic Regression, balanced accuracy 0.764062 &
    На задаче с большим числом эпох и более выраженным ERP-сигналом пайплайн
    даёт устойчивый результат заметно выше случайного уровня 0.5. \\
    \hline
  \end{tabularx}
\end{table}

Контрольные эксперименты поддерживают основной вывод данной главы. Низкие
значения на задаче реконструкции случайного визуального образа не объясняются
тем, что программный пайплайн не способен обучать модели или воспроизводимо
фиксировать результат на EEG-данных. На внешней P300-задаче тот же
инженерный контур даёт сильный перенос на удержанных испытуемых. Поэтому
наиболее вероятное ограничение основной реконструкционной постановки связано
с сочетанием малого числа примеров, слабого и вариативного сигнала визуального
воображения и высокой размерности выхода из 36 бинарных пикселей.

\section{Ограничения экспериментальной части}

Основные ограничения связаны с объёмом данных и постановкой задачи. Корпус
содержит 180 строк случайного визуального воображения, а тестовая часть
межсубъектного протокола -- 39 строк от 7 испытуемых. Для двунаправленного
межпробного протокола доступны только 27
испытуемых с двумя пробами. Это ограничивает мощность статистических проверок
и делает небольшие различия между моделями нестабильными.

Дополнительным ограничением является агрегация всей 15-секундной эпохи в одно
признаковое представление для классических моделей и один crop-вход для
нейросетевых моделей. Такая агрегация уменьшает риск завышения объёма выборки
за счёт зависимых окон, но может сглаживать кратковременные изменения
нейронной активности внутри фазы воспоминания.

Проведённая внешняя проверка на BNCI2014-001 и BNCI2014-009 подтверждает
работоспособность программного пайплайна на независимых EEG-задачах, но не
устраняет ограничение основной реконструкционной постановки: в этих корпусах
другие целевые переменные и другой экспериментальный протокол. Поэтому она
рассматривается как инженерная sanity-check валидация, а не как доказательство
декодируемости случайного визуального образа $6 \times 6$ на внешнем корпусе.
Нейросетевые модели основной задачи обучаются на малом числе строк, поэтому
даже при строгих нормировках, настраиваемых только по обучающей части, и
валидации с группировкой наблюдений их результаты следует рассматривать как
проверку реализуемости строгого протокола оценки, а не как окончательное
доказательство декодируемости случайного визуального образа из EEG.

Также не исследуются персональная адаптация, перенос обучения между
субъектами, использование EOG как дополнительного входа, генеративные модели
и реконструкция более сложных изображений. Эти направления остаются
предметом дальнейшей работы.
